% Generated draft: requires mathematical and source review.
\documentclass[11pt]{article}
\usepackage[margin=25mm]{geometry}
\usepackage{fontspec}
\setmainfont{Latin Modern Roman}
\usepackage{amsmath,amssymb,mathtools,unicode-math}
\setmathfont{Latin Modern Math}
\usepackage{enumitem,needspace,xurl,hyperref}
\setlength{\parindent}{0pt}
\setlength{\parskip}{5pt}
\setlength{\emergencystretch}{3em}
\newcommand{\N}{\mathbb{N}}
\newcommand{\Z}{\mathbb{Z}}
\newcommand{\Q}{\mathbb{Q}}
\newcommand{\R}{\mathbb{R}}
\newcommand{\C}{\mathbb{C}}
\newcommand{\F}{\mathbb{F}}
\newcommand{\A}{\mathbb{A}}
\newcommand{\PP}{\mathbb{P}}
\newcommand{\E}{\mathbb{E}}
\newcommand{\eps}{\varepsilon}
\newcommand{\dd}{\,\mathrm{d}}
\begin{document}
\section*{UnsolvedMath ID: 30001772}
\textbf{Source problem code:} catalogue entry
\subsection{Definitions and mathematical statement}
For a prime \(p\), a partition \(\lambda=(\lambda_1,\ldots,\lambda_\ell)\) of an integer \(n\) is called \(p\)-regular if no positive part occurs \(p\) or more times. Its Young diagram consists of the nodes \((r,c)\) with \(1\le r\le\ell\) and \(1\le c\le\lambda_r\). A node is removable if deleting it leaves the Young diagram of a partition. These are the usual meanings relevant to the source's phrase ``\(p\)-regular Specht module,'' which appears to refer to a Specht module labelled by a \(p\)-regular partition; that interpretation requires confirmation.

In the usual modular setting, \(S^\lambda\) denotes a Specht module for the symmetric group \(\mathfrak S_n\) over a field of characteristic \(p\), and restriction is to \(\mathfrak S_{n-1}\). The supplied record does not specify the field or explicitly identify this restriction. More importantly, it does not define the objects \(J_i^\lambda\), their construction from the restriction filtration, the ordering of removable nodes, or the precise convention for normal nodes in Kleshchev's sense.

Retaining that unresolved notation, let \(A_i\) denote the source's \(i\)-th removable node. The supplied question asks whether
\[
J_i^\lambda\ne J_{i+1}^\lambda
\quad\Longleftrightarrow\quad
A_i\text{ is normal in Kleshchev's sense},
\]
and asks for the minimal \(n\) for which this equivalence holds. The admissible indices, endpoint convention, quantification over \(p\) and \(\lambda\), and meaning of ``minimal \(n\)'' are not supplied. In particular, the record does not distinguish a least example from a least rank at which a universal assertion holds or a threshold beyond which it always holds. Consequently, this is a source-faithful statement of the proposed question with unresolved qualifications, not a fully specified conjecture.
\subsection{Short English statement}
The source asks whether jumps in a Specht restriction filtration correspond exactly to normal removable nodes, and asks for a minimal rank \(n\). The filtration notation, node conventions, hypotheses, and meaning of the minimum require clarification before a precise quantified statement can be given.
\subsection{Sources}
\begin{itemize}
\item Supplied local UnsolvedMath registry record, permanent ID \(30001772\), title ``Normal Nodes and Specht Restriction Filtrations,'' source problem code \texttt{catalogue entry}, in the Oberwolfach Reports Open Problems collection. No individual proposer is identified; the supplied proposed year is 2011.
\item Supplied report citation: \emph{Mini-Workshop: Modular Representations of Symmetric Groups and Related Objects} (2011), Oberwolfach Reports, pp.\ 1177--1205, DOI \texttt{10.4171/owr/2011/21}. Report link, also supplied under the literature-source label ``original OWR report'': \url{https://doi.org/10.4171/owr/2011/21}. The report's contents have not been independently inspected for this draft.
\item The supplied literature assessment was checked on \texttt{2026-08-21}. It labels the problem open and reports that the filtration-jump characterization and minimal-\(n\) question are posed as open problems, with no complete resolving theorem located. This is supplied dated curation derived from the OWR record, not an independently verified literature assessment; its status and cited claims require independent verification.
\end{itemize}
\end{document}
